% Esquisse d'un programme (cote 119, copy 2): archive cote 119, site page 135 (PDF page 136), typescript page 57.
% typescript. Draft compared: esquisse/pages/esq_119_059.tex (copy 1, site 59).
% AI-transcribed by Claude Opus 5.5 (claude-opus-5-5, adaptive thinking, effort high), Message Batches API, batch esq119_full (msgbatch_015xqDHkTumDysRF1qngZhRJ), 2026-11-13; not yet checked by eye.
\scanpage{119/135}
\mspage{57}

Dans le cas des espaces modérés stratifiés proprement dits (qui
ne sont pas des multiplicités à proprement parler) on peut expri-%
mer de façon commode le ``déployement'' de cette structure, i.e.
le système des espaces $X_{i_o,\ldots,i_n}$ , en termes de l'espaces mo-%
déré $X_*$ somme des $X_i$ , qui est munie d'une \uline{structure d'objet}
\uline{ordonné} (dans la catégorie des espaces modérés) ayant comme gra-%
phe $X_{**}$ de la relation d'ordre la somme des $X_{i,j}$ et des $X_i$
(ces derniers constituant la diagonale). Parmi les propriétés
essentielles de cette structure ordonnée, relevons seulement
ici que $pr_1 : X_{**} \longrightarrow X_*$ est une fibration (localement trivia-%
le) propre, et $pr_2 : X_{**} \longrightarrow X_*$ est une plongement ``bordant''.
On a une interprétation analogue du déployement d'une multipli-%
cité modérée stratifiée, en termes d'une \uline{structure de catégorie}
(remplaçant une simple structure ordonnée) ``au sens multipli-%
cités modérées'', dont l'application de composition est donnée
par les morphismes $X_{ijk} \longrightarrow X_{ik}$ ci-dessus.

\margin{135}
